%------------------------------------------------------------------------------%
\documentclass[fleqn,utf8,aspectratio=169,12pt]{beamer}

\usepackage{calculo_varias_variaveis-1}

\title{Limites de Funções de Várias Variáveis}

%------------------------------------------------------------------------------%
\begin{document}

\addtitle

%------------------------------------------------------------------------------%
\section{Limites Iterados}

%------------------------------------------------------------------------------%
\begin{question}
  Considere a função
  \[
    f(x, y) = \frac{x^2 - y^2}{x^2 + y^2}
  \]
  Calcule os limites
  \quad
  \(
    \lim_{x \to 0} \left(\lim_{y \to 0} f(x,y)\right)
  \)
  \quad
  e
  \quad
  \(
    \lim_{y \to 0} \left(\lim_{x \to 0} f(x,y)\right)
  \)
\end{question}

%------------------------------------------------------------------------------%
\begin{resolution}{Limite em $y$ depois em $x$}
  \[
    L_1
    = \lim_{x \to 0} \left( \lim_{y \to 0} \frac{x^2 - y^2}{x^2 + y^2} \right)
  \]

  \pause
  Para $x$ fixo diferente de zero
  \[
    \lim_{y \to 0} \frac{x^2 - y^2}{x^2 + y^2}
    \pause
    \;=\; \frac{x^2}{x^2}
    \pause
    \;=\; 1
  \]

  \pause
  Logo
  \[
    L_1
    \pause
    \;=\; \lim_{x \to 0} \left(\lim_{y \to 0} f(x,y)\right)
    \pause
    \;=\; \lim_{x \to 0} 1
    \pause
    \;=\; 1
  \]

  ~
\end{resolution}

%------------------------------------------------------------------------------%
\begin{resolution}{Limite em $x$ depois em $y$}
  \[
    L_2
    = \lim_{y \to 0} \left( \lim_{x \to 0} \frac{x^2 - y^2}{x^2 + y^2} \right)
  \]

  \pause
  Para $y$ diferente de zero
  \[
    \lim_{x \to 0} \frac{x^2 - y^2}{x^2 + y^2}
    \pause
    \;=\; \frac{-y^2}{y^2}
    \pause
    \;=\; -1
  \]

  \pause
  Logo
  \[
    L_2
    \pause
    \;=\; \lim_{y \to 0} \left(\lim_{x \to 0} f(x,y)\right)
    \pause
    \;=\; \lim_{y \to 0} (-1)
    \pause
    \;=\; -1
  \]

  \pause
  Não podemos definir o limite em várias variáveis dessa forma
\end{resolution}

%------------------------------------------------------------------------------%


%------------------------------------------------------------------------------%
\section{Limites}

%------------------------------------------------------------------------------%
\begin{frame}{Definição}
  Uma função
  \(f\colon R \subset \R^2 \to \R\), \;
  \pause
  \(f(x, y)\),

  \pause
  se aproxima do \emph{limite} $L$
  a medida que \((x, y)\) se aproxima de \((a, b)\)

  \pause
  se, para todo \emph{$\varepsilon>0$,}
  \pause
  existe um \emph{\(\delta(\varepsilon)>0\)}
  \pause
  tal que,
  \pause
  para todo \((x, y)\) no domínio de $f$

  \pause
  \[
    \abs{f(x, y)-L} < \varepsilon
    \quad\text{ sempre que }\quad
    0 < \sqrt{ \left(x-a\right)^2 + \left(y-b\right)^2 } < \delta(\varepsilon)
  \]

  \pause
  Obs: \((a, b)\) não precisa estar no domínio de $f$
\end{frame}

%------------------------------------------------------------------------------%
\begin{frame}{Notação}
  \[
    \lim_{(x, y)\to(a, b)} f(x, y) = L
  \]
  \pause
  \vspace{\baselineskip}
  \[
    f(x, y) \to L \quad\text{ quando }\quad (x, y)\to(a, b)
  \]
\end{frame}

%------------------------------------------------------------------------------%
\begin{question}
\begin{columns}[t]
\column{0.5\linewidth}
  Encontre os limites
  \vspace{0.5\baselineskip}
  \begin{enumerate}
    \setlength\itemsep{1.5em}
    \item \(\lim_{(x, y)\to(a, b)} k\)
    \item \(\lim_{(x, y)\to(a, b)} x\)
    \item \(\lim_{(x, y)\to(a, b)} y\)
    \item \(\lim_{(x, y)\to(a, b)} xy\)
  \end{enumerate}
\column{0.5\linewidth}
  \pause
  Respostas
  \pause
  \vspace{0.5\baselineskip}
  \begin{enumerate}[<+->]
    \setlength\itemsep{1.5em}
    \item \(\lim_{(x, y)\to(a, b)} k  = k  \)
    \item \(\lim_{(x, y)\to(a, b)} x  = a  \)
    \item \(\lim_{(x, y)\to(a, b)} y  = b  \)
    \item \(\lim_{(x, y)\to(a, b)} xy = ab \)
  \end{enumerate}
\end{columns}
\end{question}

%------------------------------------------------------------------------------%


%------------------------------------------------------------------------------%
\section{Propriedades}

%------------------------------------------------------------------------------%
\begin{frame}{Teorema -- Unicidade do Limite}
  Se uma função \(f(x, y)\) possui um limite $L$

  quando \((x, y)\) se aproxima de \((a, b)\),

  \pause
  então \emph{$L$ é único}
\end{frame}

%------------------------------------------------------------------------------%
\begin{frame}{Teorema -- Propriedades dos limites de várias variáveis}
  As propriedades a seguir são verdadeiras se
  \begin{itemize}[<+->]
    \setlength\itemsep{1.5em}
    \item $L$, $M$ e $k$ forem números reais
    \item \(\lim_{(x, y)\to(a, b)} f(x, y) = L\)
    \item \(\lim_{(x, y)\to(a, b)} g(x, y) = M\)
  \end{itemize}
\end{frame}

%------------------------------------------------------------------------------%
\begin{frame}{Regra da Soma e Diferença}
  \[
    \lim_{(x, y)\to(a, b)} \left[f(x, y) + g(x, y)\right]
    \pause =
    \lim_{(x, y)\to(a, b)} f(x, y) + \lim_{(x, y)\to(a, b)} g(x, y)
    \pause =
    L + M
  \]
  \pause
  \vspace{\baselineskip}
  \[
    \lim_{(x, y)\to(a, b)} \left[f(x, y) - g(x, y)\right]
    \pause =
    \lim_{(x, y)\to(a, b)} f(x, y) - \lim_{(x, y)\to(a, b)} g(x, y)
    \pause =
    L - M
  \]
\end{frame}

%------------------------------------------------------------------------------%
\begin{frame}{Regra da Multiplicação por Constante e do Produto}
  \[
    \lim_{(x, y)\to(a, b)} \left[\,k\, f(x, y)\right]
    \pause =
    k \lim_{(x, y)\to(a, b)} f(x, y)
    \pause =
    k L
  \]
  \pause
  \vspace{\baselineskip}
  \[
    \lim_{(x, y)\to(a, b)} \left[f(x, y) g(x, y)\right]
    \pause =
    \left[\lim_{(x, y)\to(a, b)} f(x, y) \right]
    \left[\lim_{(x, y)\to(a, b)} g(x, y) \right]
    \pause =
    LM
  \]
\end{frame}

%------------------------------------------------------------------------------%
\begin{frame}{Regra do Quociente (Divisão)}
  \[
    \lim_{(x, y)\to(a, b)} \left[ \,\frac{\,f(x, y)\,}{g(x, y)} \, \right]
    \pause =
    \dfrac{\displaystyle\lim_{(x, y)\to(a, b)} f(x, y)\,}
          {\displaystyle\lim_{(x, y)\to(a, b)} g(x, y)\,}
    \pause =
    \frac{L}{M}
  \]

  \pause
  \emph{Desde que \(M\neq 0\)}
\end{frame}

%------------------------------------------------------------------------------%
\begin{frame}{Regra da Potência e da Raiz}
  Para $n$ inteiro positivo
  \pause
  \[
    \lim_{(x, y)\to(a, b)} \left[f(x, y)\right]^n
    \pause =
    \left[\lim_{(x, y)\to(a, b)} f(x, y)\right]^n
    \pause =
    L^n
  \]
  \pause
  \vspace{\baselineskip}
  \[
    \lim_{(x, y)\to(a, b)} \sqrt[n]{f(x, y)}
    \pause =
    \sqrt[n]{\lim_{(x, y)\to(a, b)} f(x, y)}
    \pause =
    \sqrt[n]{L}
  \]

  \pause
  Se $n$ for par $L$ precisa ser positivo
\end{frame}

%------------------------------------------------------------------------------%
\begin{frame}{Teorema do Confronto}
  Se
  \[
    g(x, y) \leq f(x, y) \leq h(x, y)
  \]
  \pause
  para todo \((x, y) \neq (a, b)\),
  \pause
  em um disco centrado em \((a, b)\)
  \pause
  e
  \[
    \lim_{(x, y)\to(a, b)} g(x, y)
    \pause
    \;=\;
    \lim_{(x, y)\to(a, b)} h(x, y)
    \pause
    \;=\;
    L
  \]
  \pause
  então
  \[
    \lim_{(x, y)\to(a, b)} f(x, y) \;=\; L
  \]
\end{frame}

%------------------------------------------------------------------------------%
\section{Exemplos}

%------------------------------------------------------------------------------%
\begin{question}
  Calcule o limite
  \(
    \lim_{(x, y)\to(0, 1)} \frac{x-xy+3}{x^2y +5xy - y^3}
  \),
  se ele existir

  \pause
  \vspace{2\baselineskip}
  Assumindo que todos os limites necessários para aplicar as propriedades existam

  (Isso pode não ser verdade)
\end{question}

%------------------------------------------------------------------------------%
\begin{resolution}{Solução}
\begin{align*}
  \onslide<+->{\lim_{(x, y)\to(0, 1)} \frac{x-xy+3}{x^2y +5xy - y^3}}
  \onslide<+->{& = \frac{\displaystyle \lim_{(x, y)\to(0, 1)}(x-xy+3)}
           {\displaystyle \lim_{(x, y)\to(0, 1)}(x^2y +5xy - y^3)} \\}
  \onslide<+->{& = \frac{\displaystyle \lim_{(x, y)\to(0, 1)} x -
                          \lim_{(x, y)\to(0, 1)} xy +
                          \lim_{(x, y)\to(0, 1)} 3}
           {\displaystyle \lim_{(x, y)\to(0, 1)} x^2y +
                          \lim_{(x, y)\to(0, 1)} 5xy -
                          \lim_{(x, y)\to(0, 1)} y^3} \\}
  \onslide<+->{& = \frac{0 - 0\times 1 + 3}{0^2\times 1 + 5\times 0 \times 1 - 1^3} \\}
  \onslide<+->{& = \frac{3}{-1}}
  \onslide<+->{= -3}
\end{align*}
\end{resolution}

%------------------------------------------------------------------------------%

%------------------------------------------------------------------------------%
\begin{question}
  Calcule o limite
  \(
    \lim_{(x, y)\to(0, 0)} \frac{x^2-xy}{\,\sqrt{x}-\sqrt{y}\,}
  \)
  se ele existir

  \pause
  \vspace{\baselineskip}
  Assumindo que todos os limites necessários para aplicar as propriedades existam

  (Isso pode não ser verdade)
\end{question}

%------------------------------------------------------------------------------%
\begin{resolution}{Solução}
\begin{align*}
  \onslide<+->{\lim_{(x, y)\to(0, 0)} \frac{x^2-xy}{\,\sqrt{x}-\sqrt{y}\,}}
  \onslide<+->{& = \frac{\displaystyle \lim_{(x, y)\to(0, 0)} \left(x^2-xy\right)}
           {\displaystyle \lim_{(x, y)\to(0, 0)} \left(\sqrt{x}-\sqrt{y}\right)} \\[2mm]}
  \onslide<+->{& = \frac{\displaystyle \lim_{(x, y)\to(0, 0)} x^2 - \lim_{(x, y)\to(0, 0)}xy}
           {\displaystyle \lim_{(x, y)\to(0, 0)} \sqrt{x} - \lim_{(x, y)\to(0, 0)}\sqrt{y}} \\[2mm]}
  \onslide<+->{& =   \frac{0-0}{0-0} }
  \onslide<+->{\;=\; \frac{0}{0} }
  \onslide<+->{\;=\; 1}
\end{align*}

\onslide<+->{\textcolor{red}{\textbf{Não existe divisão por zero !!!}}}

~

~
\end{resolution}

%------------------------------------------------------------------------------%
\begin{resolution}{Solução}
\begin{align*}
  \,\lim_{(x, y)\to(0, 0)} \frac{x^2-xy}{\,\sqrt{x}-\sqrt{y}\,}
  & = \frac{\displaystyle \lim_{(x, y)\to(0, 0)} \left(x^2-xy\right)}
    {\displaystyle \lim_{(x, y)\to(0, 0)} \left(\sqrt{x}-\sqrt{y}\right)} \\[2mm]
  & = \frac{\displaystyle \lim_{(x, y)\to(0, 0)} x^2 - \lim_{(x, y)\to(0, 0)}xy}
    {\displaystyle \lim_{(x, y)\to(0, 0)} \sqrt{x} - \lim_{(x, y)\to(0, 0)}\sqrt{y}} \\[2mm]
  & =  \frac{0-0}{0-0}
  \;=\; \frac{0}{0}
  \qquad\nexists
\end{align*}

\emph{Então o limite não existe?}

~

~
\end{resolution}

%------------------------------------------------------------------------------%
\begin{resolution}{Solução}
  \begin{align*}
    \,\lim_{(x, y)\to(0, 0)} \frac{x^2-xy}{\,\sqrt{x}-\sqrt{y}\,}
    & \color{red}
    = \frac{\displaystyle \lim_{(x, y)\to(0, 0)} \left(x^2-xy\right)}
    {\displaystyle \lim_{(x, y)\to(0, 0)} \left(\sqrt{x}-\sqrt{y}\right)} \\[2mm]
    & \color{red} = \frac{\displaystyle \lim_{(x, y)\to(0, 0)} x^2 - \lim_{(x, y)\to(0, 0)}xy}
    {\displaystyle \lim_{(x, y)\to(0, 0)} \sqrt{x} - \lim_{(x, y)\to(0, 0)}\sqrt{y}} \\[2mm]
    & \color{red} = \frac{0-0}{0-0}
    \;=\; \frac{0}{0}
  \end{align*}

  Não podemos aplicar a propriedade do quociente

  \pause
  \emph{Precisamos remover a indeterminação}

  ~
\end{resolution}

%------------------------------------------------------------------------------%
\begin{resolution}{Solução}
  \onslide<+->{Ao calcular o limite consideramos apenas os pontos dentro do domínio da função}

  \onslide<+->{Para \(x\neq y\)}

  \begin{align*}
    \onslide<+->{\frac{x^2-xy}{\,\sqrt{x}-\sqrt{y}\,}}
    \onslide<+->{& = \frac{\left(x^2-xy\right) \emph{\left(\sqrt{x}+\sqrt{y}\right)}}
                          {\left(\sqrt{x}-\sqrt{y}\right)\emph{\left(\sqrt{x}+\sqrt{y}\right)}} \\[1mm]}
    \onslide<+->{& = \frac{x\left(\emph{x-y}\right)\left(\sqrt{x}+\sqrt{y}\right)}{\emph{x-y}} \\[1mm]}
    \onslide<+->{& = x\left(\sqrt{x}+\sqrt{y}\right)}
  \end{align*}
\end{resolution}

%------------------------------------------------------------------------------%
\begin{resolution}{Solução}
\onslide<+->{Calculando o limite}

\begin{align*}
  \onslide<+->{\lim_{(x, y)\to(0, 0)} \frac{x^2-xy}{\,\sqrt{x}-\sqrt{y}\,}}
  \onslide<+->{& = \lim_{(x, y)\to(0, 0)} x\left(\sqrt{x}+\sqrt{y}\right) \\[1mm]}
  \onslide<+->{& = \left(\lim_{(x, y)\to(0, 0)} x\right)
    \left[
    \left(\lim_{(x, y)\to(0, 0)}\sqrt{x}\right) +
    \left(\lim_{(x, y)\to(0, 0)}\sqrt{y}\right)\right] \\[2mm]}
  \onslide<+->{& = 0\left(\sqrt{0}+\sqrt{0}\right) \\[2mm]}
  \onslide<+->{& = 0}
\end{align*}
\end{resolution}

%------------------------------------------------------------------------------%

%------------------------------------------------------------------------------%
\begin{question}
  Calcule o limite
  \(
    \lim_{(x,y) \to (4, 3)} \frac{\sqrt{x}-\sqrt{y+1}}{x-y-1}
  \)
\end{question}

%------------------------------------------------------------------------------%
\begin{resolution}{Manipulando a Função}
  \onslide<+->{No interior do domínio da função}

  \onslide<+->{Multiplicando pelo conjugado}

  \begin{align*}
    \onslide<+->{\frac{\sqrt{x}-\sqrt{y+1}}{x-y-1}}
    \onslide<+->{& = \frac{\sqrt{x} - \sqrt{y+1}}{x - y - 1}
        \times
        \frac{\sqrt{x} + \sqrt{y+1}}{\sqrt{x} + \sqrt{y+1}} \\}
    \onslide<+->{& = \frac{x - (y+1)}{(x - y - 1)\left(\sqrt{x} + \sqrt{y+1}\right)} \\}
    \onslide<+->{& = \frac{1}{\sqrt{x} + \sqrt{y+1}}}
  \end{align*}
\end{resolution}

%------------------------------------------------------------------------------%
\begin{resolution}{Calculando o Limite}
  \begin{align*}
    \onslide<+->{\lim_{(x,y) \to (4, 3)} \frac{\sqrt{x}-\sqrt{y+1}}{x-y-1}}
    \onslide<+->{& = \lim_{(x,y) \to (4, 3)} \frac{1}{\sqrt{x} + \sqrt{y+1}} \\}
    \onslide<+->{& = \frac{1}{\sqrt{4} + \sqrt{3+1}} \\}
    \onslide<+->{& = \frac{1}{2+2} \\}
    \onslide<+->{& = \frac{1}{4}}
  \end{align*}
\end{resolution}

%------------------------------------------------------------------------------%

%------------------------------------------------------------------------------%
\begin{question}
  Calcule o limite
  \(
    \lim_{(x,y) \to (0,1)} \frac{\ln(y)\sin(x)}{x}
  \)
\end{question}

%------------------------------------------------------------------------------%
\begin{resolution}{Solução}
  \onslide<+->{Limite fundamental
  \[
    \lim_{x\to 0} \frac{\sin(x)}{x} = 1
  \]}

  \begin{align*}
    \onslide<+->{\lim_{(x,y) \to (0,1)} \frac{\ln(y)\sin(x)}{x}}
    \onslide<+->{& = \left(\lim_{(x,y) \to (0,1)} \ln(y)\right)
                     \left(\lim_{(x,y) \to (0,1)} \frac{\sin(x)}{x}\right) \\}
    \onslide<+->{& = \ln(1) \times 1 \\}
    \onslide<+->{& = 0}
  \end{align*}
\end{resolution}

%------------------------------------------------------------------------------%

%------------------------------------------------------------------------------%
\begin{question}
  Calcule o limite
  \(
    \lim_{(x,y) \to (0,0)} \frac{\sqrt{x^2 + y^2 + 1} - 1}{x^2 + y^2}
  \)
\end{question}

%------------------------------------------------------------------------------%
\begin{resolution}{Multiplicando pelo conjugado}
  \onslide<+->{No domínio da função}
  \begin{align*}
    \onslide<+->{\frac{\sqrt{x^2 + y^2 + 1} - 1}{x^2 + y^2}}
    \onslide<+->{& = \frac{\sqrt{x^2 + y^2 + 1} - 1}{x^2 + y^2}
                     \times
                     \frac{\sqrt{x^2 + y^2 + 1} + 1}{\sqrt{x^2 + y^2 + 1} + 1} \\}
    \onslide<+->{& = \frac{x^2 + y^2}{(x^2 + y^2)(\sqrt{x^2 + y^2 + 1} + 1)} \\}
    \onslide<+->{& = \frac{1}{\sqrt{x^2 + y^2 + 1} + 1}}
  \end{align*}
\end{resolution}

%------------------------------------------------------------------------------%
\begin{resolution}{Calculando o Limite}
  \begin{align*}
    \onslide<+->{\lim_{(x,y) \to (0,0)} \frac{\sqrt{x^2 + y^2 + 1} - 1}{x^2 + y^2}}
    \onslide<+->{& = \lim_{(x,y) \to (0,0)} \frac{1}{\sqrt{x^2 + y^2 + 1} + 1}  \\}
    \onslide<+->{& = \frac{1}{\sqrt{0^2 + 0^2 + 1} + 1}   \\}
    \onslide<+->{& = \frac{1}{2}}
  \end{align*}
\end{resolution}

%------------------------------------------------------------------------------%


%------------------------------------------------------------------------------%
\section{Lista Mínima}

%------------------------------------------------------------------------------%
\begin{frame}{Lista Mínima}
  Cálculo Vol. 2 do Thomas 12\textsuperscript{a} ed. -- Seção 14.2
  \begin{enumerate}
    \item Estudar o texto da seção
    \item Resolver os exercícios: 1, 5, 7, 9, 11, 15, 17, 23, 27, 55, 57
  \end{enumerate}

  \vfill
  Atenção: A prova é baseada no livro, não nas apresentações
\end{frame}

%------------------------------------------------------------------------------%
\end{document}
%------------------------------------------------------------------------------%
